% !TEX root = ../main.tex
\begin{ouverture}[colorimetrie]{5}{Titrages colorimétriques}{Partie 1 · Constitution et transformations de la matière}
{Titrage\enspace·\enspace Réaction support\enspace·\enspace Équivalence\enspace·\enspace Volume équivalent\enspace·\enspace Changement de teinte}
\objectif{Décrire le principe d'un titrage direct et les qualités de sa réaction support.}
\objectif{Relier l'équivalence au changement de réactif limitant et aux proportions stœchiométriques.}
\objectif{Établir la relation entre les quantités de matière de réactifs à l'équivalence et en déduire une concentration.}
\objectif{Expliquer ou prévoir le changement de teinte observé à l'équivalence.}
\objectif{Réaliser un titrage avec repérage colorimétrique et exploiter le volume équivalent.}
\prerequis{Couples oxydant-réducteur et équations d'oxydoréduction (chapitre~3).}
\prerequis{Avancement, réactif limitant, mélange stœchiométrique (chapitre~4).}
\prerequis{Relation $n=c\times V$ et dilution (chapitre~1).}
\end{ouverture}
\essentiel
\essSection{1}{Le principe d'un titrage}
\begin{definition}[Titrage direct]
\textbf{Titrer} une espèce chimique dans une solution, c'est déterminer sa quantité de matière, ou sa concentration, à
l'aide d'une réaction chimique avec une autre espèce de concentration connue, le \textbf{réactif titrant}. Cette
\textbf{réaction support du titrage} doit être \textbf{totale}, \textbf{rapide} et \textbf{unique} : le réactif titrant ne doit réagir
qu'avec l'espèce titrée. En classe de première, les réactions support sont des réactions d'oxydoréduction.
\end{definition}
\schema{05-montage}
\essSection{2}{L'équivalence}
\begin{definition}[Équivalence d'un titrage]
L'\textbf{équivalence} est atteinte lorsque les réactifs ont été introduits dans les \textbf{proportions stœchiométriques}. Elle
correspond à un \textbf{changement de réactif limitant} :
\begin{questions}
\item avant l'équivalence, le réactif titrant versé est limitant : il est entièrement consommé dès son ajout ;
\item après l'équivalence, l'espèce titrée a disparu : le réactif titrant versé s'accumule dans le mélange.
\end{questions}
Le volume de solution titrante versé à l'équivalence est le \textbf{volume équivalent} $V_\mathrm{E}$.
\end{definition}
\schema{05-evolution}
\essSection{3}{Exploiter l'équivalence}
\begin{propriete}[Relation à l'équivalence]
Pour la réaction support $a\,\mathrm{A}+b\,\mathrm{B}\longrightarrow\dots$, où A est l'espèce titrée et B le réactif titrant de concentration $c_\mathrm{B}$,
les proportions sont stœchiométriques à l'équivalence :
\[\frac{n_0(\mathrm{A})}{a}=\frac{n_\mathrm{E}(\mathrm{B})}{b}=\frac{c_\mathrm{B}\,V_\mathrm{E}}{b}.\]
\end{propriete}
\begin{savoirfaire}[Exploiter un titrage]
\begin{questions}
\item Écrire l'équation de la réaction support, à partir des couples si nécessaire.
\item Calculer la quantité de réactif titrant versé à l'équivalence : $n_\mathrm{E}(\mathrm{B})=c_\mathrm{B}\times V_\mathrm{E}$.
\item Écrire la relation à l'équivalence et en déduire $n_0(\mathrm{A})$, puis la concentration $c_\mathrm{A}=\dfrac{n_0(\mathrm{A})}{V_\mathrm{A}}$.
\item Tenir compte, s'il y a lieu, de la dilution ou du prélèvement effectués avant le titrage.
\end{questions}
\end{savoirfaire}
\begin{exemple}
On titre $V_1=\qty{10,0}{mL}$ d'une solution d'ions fer(II) par une solution de permanganate de potassium à
$c_2=\qty{2,00e-2}{mol.L^{-1}}$ ; $V_\mathrm{E}=\qty{12,0}{mL}$. Avec $\ce{MnO4- + 8H+ + 5Fe^{2+} -> Mn^{2+} + 5Fe^{3+} + 4H2O}$ :
\[\frac{n_0(\ce{Fe^{2+}})}{5}=c_2\,V_\mathrm{E}\qquad\text{donc}\qquad n_0(\ce{Fe^{2+}})=5\times\num{2.00e-2}\times\num{12.0e-3}=\qty{1,20e-3}{mol}\]
et $c_1=\dfrac{\num{1.20e-3}}{\num{10.0e-3}}=\qty{0,120}{mol.L^{-1}}$.
\end{exemple}
\essSection{4}{Repérer l'équivalence par un changement de teinte}
\begin{propriete}[Repérage colorimétrique]
Si une espèce colorée intervient, la teinte du mélange change à l'équivalence :
\begin{questions}
\item si le réactif titrant est coloré, comme l'ion permanganate violet, sa teinte \textbf{persiste} à partir de l'équivalence ;
\item si l'espèce titrée est colorée, comme le diiode jaune-brun, sa teinte \textbf{disparaît} à l'équivalence ;
\item sinon, on ajoute un indicateur coloré : l'empois d'amidon colore en bleu intense une solution contenant du diiode.
\end{questions}
\end{propriete}
\schema{05-couleurs}
\begin{python}
c2 = 2.00e-2      # mol/L : ions permanganate
n1 = 1.20e-3      # mol : ions fer(II)
for V in [0, 4, 8, 12, 16]:          # mL
    n2 = c2 * V * 1e-3
    x = min(n2 / 1, n1 / 5)
    nFe = n1 - 5 * x
    nMn = n2 - x
    print(f"V = {V:2d} mL : Fe2+ {nFe * 1e3:.2f} mmol ; MnO4- {nMn * 1e3:.2f} mmol")
\end{python}
\begin{console}
V =  0 mL : Fe2+ 1.20 mmol ; MnO4- 0.00 mmol
V =  4 mL : Fe2+ 0.80 mmol ; MnO4- 0.00 mmol
V =  8 mL : Fe2+ 0.40 mmol ; MnO4- 0.00 mmol
V = 12 mL : Fe2+ 0.00 mmol ; MnO4- 0.00 mmol
V = 16 mL : Fe2+ 0.00 mmol ; MnO4- 0.08 mmol
\end{console}
\noindent\capacite{Capacité numérique} Pour chaque volume versé, le programme calcule l'avancement maximal \code{x} comme au chapitre~4, puis
les quantités restantes d'ions fer(II) et permanganate. Avant \qty{12}{mL}, il ne reste jamais d'ions permanganate : ils sont
limitants. Après, il ne reste plus d'ions fer(II) et les ions permanganate s'accumulent : la teinte violette apparaît.
\begin{attention}
L'équivalence n'est pas la fin de la réaction : c'est l'instant où les réactifs ont été apportés dans les proportions
stœchiométriques. L'eau distillée ajoutée dans l'erlenmeyer pour rincer ses parois ne change pas $n_0(\mathrm{A})$, donc ne
change pas $V_\mathrm{E}$. Près de l'équivalence, on verse la solution titrante goutte à goutte.
\end{attention}
\begin{experience}[Titrer la vitamine C d'un comprimé]
Un comprimé « vitamine C 500 » est dissous dans une fiole jaugée de \qty{200,0}{mL}. On prélève \qty{10,0}{mL} de cette solution, on
ajoute quelques gouttes d'empois d'amidon, puis on verse une solution de diiode à \qty{1,00e-2}{mol.L^{-1}}. Le diiode est réduit
dès son ajout, tant qu'il reste de la vitamine C ; la teinte bleue persiste à partir de $V_\mathrm{E}=\qty{14,2}{mL}$. On en déduit
\qty{2,84e-3}{mol} de vitamine C dans le comprimé, soit \qty{500}{mg} : l'indication est vérifiée.
\end{experience}
\begin{rediger}
Écrire l'équation de la réaction support, définir l'équivalence, écrire la relation à l'équivalence avec les nombres
stœchiométriques, puis calculer en gardant les unités. Ne pas oublier le facteur de dilution ou de prélèvement, et
comparer le résultat à la valeur attendue.
\end{rediger}
\exercices

\rubrique{Découverte}
\begin{blocLivre}{ex-5-1}
\begin{exercice}[1]{Le vocabulaire du titrage}
On titre une solution jaune-orangé de diiode par une solution incolore de thiosulfate de sodium.
\schema{05-ex01}
\begin{questions}
\item Nommer les éléments numérotés de 1 à 5.
\item Quelle solution est placée dans la burette ? dans l'erlenmeyer ? Avec quelle verrerie prélève-t-on la solution titrée ?
\item Pourquoi agite-t-on le mélange ? Pourquoi place-t-on une feuille blanche sous l'erlenmeyer ?
\item Définir l'équivalence du titrage.
\end{questions}
\end{exercice}
\end{blocLivre}
\begin{blocLivre}{ex-5-2}
\begin{exercice}[1]{Titrer une solution de diiode}
On titre $V_1=\qty{20,0}{mL}$ d'une solution de diiode par une solution de thiosulfate de sodium à $c_2=\qty{0,100}{mol.L^{-1}}$. La
réaction support a pour équation $\ce{I2 + 2S2O3^{2-} -> 2I- + S4O6^{2-}}$. L'équivalence est obtenue pour $V_\mathrm{E}=\qty{12,4}{mL}$.
\begin{questions}
\item Quel est le réactif titrant ? Quelle est l'espèce titrée ?
\item Écrire la relation entre les quantités de matière à l'équivalence.
\item Calculer la quantité de diiode titrée, puis sa concentration $c_1$.
\item Seul le diiode est coloré. Comment l'équivalence est-elle repérée ?
\end{questions}
\end{exercice}
\end{blocLivre}
\begin{blocLivre}{ex-5-3}
\begin{exercice}[1]{Avant, à et après l'équivalence}
On titre une solution d'ions fer(II), presque incolore, par une solution violette d'ions permanganate, en milieu acide.
On s'intéresse au contenu de l'erlenmeyer pour trois volumes versés : $V<V_\mathrm{E}$, $V=V_\mathrm{E}$ et $V>V_\mathrm{E}$.
\begin{questions}
\item Pour $V<V_\mathrm{E}$, quel est le réactif limitant ? Quelles espèces réactives restent dans l'erlenmeyer ?
\item Que peut-on dire des quantités de réactifs à l'équivalence ?
\item Pour $V>V_\mathrm{E}$, quel est le réactif limitant ? Quelle espèce s'accumule ?
\item Décrire la teinte du mélange dans chacun des trois cas.
\end{questions}
\end{exercice}
\end{blocLivre}
\begin{blocLivre}{ex-5-4}
\begin{exercice}[1]{Prévoir le changement de teinte}
Le tableau décrit quatre titrages. Prévoir, pour chacun, le changement de teinte observé à l'équivalence.
\par\smallskip\noindent\hfill\begin{tabular}{clll}
\toprule
 & Espèce titrée & Réactif titrant & Remarque \\
\midrule
a & $\ce{Fe^{2+}}$, vert très pâle & $\ce{MnO4-}$, violet & produits presque incolores \\
b & $\ce{I2}$, jaune-brun & $\ce{S2O3^{2-}}$, incolore & produits incolores \\
c & vitamine C, incolore & $\ce{I2}$, jaune-brun & empois d'amidon ajouté \\
d & $\ce{MnO4-}$, violet & $\ce{Fe^{2+}}$, vert très pâle & produits presque incolores \\
\bottomrule
\end{tabular}\hfill\null\par\smallskip
\begin{questions}
\item Titrage a.
\item Titrage b ; comment rendre le repérage plus précis ?
\item Titrage c.
\item Titrage d.
\end{questions}
\end{exercice}
\end{blocLivre}
\begin{blocLivre}{ex-5-5}
\begin{exercice}[1]{Lire une burette et conclure}
On titre $V_1=\qty{20,0}{mL}$ d'une solution d'ions fer(II) par une solution de permanganate de potassium à \qty{2,00e-2}{mol.L^{-1}}.
Le schéma montre la burette à l'équivalence.
\schema{05-ex05}
\begin{questions}
\item Lire le volume équivalent.
\item Pourquoi lit-on le bas du ménisque, l'œil à sa hauteur ?
\item Écrire la relation à l'équivalence ; l'équation est celle de l'essentiel.
\item Calculer la concentration $c_1$ des ions fer(II).
\end{questions}
\end{exercice}
\end{blocLivre}
\rubrique{Entraînement}
\begin{blocLivre}{ex-5-6}
\begin{exercice}[2]{Simuler un titrage en Python}
On reprend le programme de l'essentiel.
\begin{questions}
\item \capacite{Capacité numérique}Que représente \code{x} ? Pourquoi est-il calculé avec la fonction \code{min} ?
\item Prévoir la ligne affichée pour $V=\qty{6}{mL}$, puis pour $V=\qty{20}{mL}$.
\item Quelle ligne correspond à l'équivalence ? Comment modifier la liste des volumes pour mieux la localiser ?
\item On veut simuler le titrage de \qty{0,400}{mmol} de diiode par du thiosulfate à \qty{5,00e-2}{mol.L^{-1}}. Écrire les lignes à modifier.
\end{questions}
\end{exercice}
\end{blocLivre}
\begin{blocLivre}{ex-5-7}
\begin{exercice}[2]{Un médicament contre l'anémie}
Un médicament contre l'anémie annonce \qty{80}{mg} de fer par comprimé, sous forme d'ions fer(II). On dissout un comprimé dans
de l'acide sulfurique dilué, dans une fiole jaugée de \qty{100,0}{mL}. On titre $V_1=\qty{10,0}{mL}$ de cette solution par du
permanganate de potassium à \qty{2,00e-3}{mol.L^{-1}} : $V_\mathrm{E}=\qty{14,3}{mL}$.
\begin{donnees}
Couples $\ce{MnO4-}/\ce{Mn^{2+}}$ et $\ce{Fe^{3+}}/\ce{Fe^{2+}}$ ; $M(\ce{Fe})=\qty{55,8}{g.mol^{-1}}$.
\end{donnees}
\begin{questions}
\item Établir l'équation de la réaction support. Pourquoi faut-il opérer en milieu acide ?
\item Calculer la quantité d'ions fer(II) dans le volume titré.
\item En déduire la masse de fer d'un comprimé et la comparer à l'indication.
\item Quel volume équivalent aurait-on obtenu pour exactement \qty{80}{mg} de fer ?
\end{questions}
\end{exercice}
\end{blocLivre}
\begin{blocLivre}{ex-5-8}
\begin{exercice}[2]{Un antiseptique iodé}
Un antiseptique annonce « 1~\% d'iode disponible », c'est-à-dire \qty{1}{g} de diiode pour \qty{100}{mL}. On le dilue dix fois, puis on
titre $V_1=\qty{10,0}{mL}$ de la solution diluée par du thiosulfate de sodium à \qty{1,00e-2}{mol.L^{-1}} : $V_\mathrm{E}=\qty{7,90}{mL}$.
\begin{donnees}
$\ce{I2 + 2S2O3^{2-} -> 2I- + S4O6^{2-}}$ ; $M(\ce{I})=\qty{126,9}{g.mol^{-1}}$.
\end{donnees}
\begin{questions}
\item Décrire le protocole de la dilution, en précisant la verrerie.
\item Comment repère-t-on l'équivalence ?
\item Calculer la concentration en diiode de la solution diluée, puis de l'antiseptique.
\item Calculer la concentration en masse de diiode de l'antiseptique et conclure.
\end{questions}
\end{exercice}
\end{blocLivre}
\begin{blocLivre}{ex-5-9}
\begin{exercice}[2]{Contrôler un détachant}
Un détachant pour le bois est vendu comme de l'acide oxalique pur, $\ce{H2C2O4}$, cristallisé avec deux molécules d'eau
($M=\qty{126,0}{g.mol^{-1}}$). On dissout \qty{1,00}{g} de poudre dans une fiole jaugée de \qty{100,0}{mL}. On titre, à chaud et en milieu
acide, $V_1=\qty{10,0}{mL}$ de cette solution par du permanganate de potassium à \qty{2,00e-2}{mol.L^{-1}} : $V_\mathrm{E}=\qty{15,6}{mL}$.
\begin{donnees}
Couples $\ce{MnO4-}/\ce{Mn^{2+}}$ et $\ce{CO2}/\ce{H2C2O4}$.
\end{donnees}
\begin{questions}
\item Écrire la demi-équation du couple $\ce{CO2}/\ce{H2C2O4}$ et établir l'équation de la réaction support.
\item Décrire le changement de teinte à l'équivalence.
\item Calculer la masse d'acide oxalique cristallisé dans l'échantillon et la pureté du produit.
\item Quel volume équivalent aurait donné un produit pur ?
\end{questions}
\end{exercice}
\end{blocLivre}
\begin{blocLivre}{ex-5-10}
\begin{exercice}[2]{Bien choisir la solution titrante}
On veut titrer $V_1=\qty{10,0}{mL}$ d'une solution d'ions fer(II) de concentration voisine de \qty{0,050}{mol.L^{-1}} avec une burette de
\qty{25}{mL}. On dispose de trois solutions de permanganate : \qty{1,00e-1}{mol.L^{-1}}, \qty{1,00e-2}{mol.L^{-1}} et \qty{5,00e-3}{mol.L^{-1}}.
\begin{questions}
\item Exprimer le volume équivalent en fonction de la concentration $c_2$ du permanganate.
\item Calculer le volume équivalent attendu avec chaque solution.
\item La lecture d'un volume sur la burette est entachée d'une incertitude d'environ \qty{0,1}{mL}. Calculer l'incertitude relative
dans chaque cas.
\item Quelle solution choisir ? Justifier.
\end{questions}
\end{exercice}
\end{blocLivre}
\begin{blocLivre}{ex-5-11}
\begin{exercice}[2]{Une eau oxygénée a-t-elle vieilli ?}
Une eau oxygénée « à 10 volumes » contient en principe \qty{0,89}{mol.L^{-1}} de peroxyde d'hydrogène (chapitre~3). On la dilue dix fois,
puis on titre $V_1=\qty{10,0}{mL}$ de la solution diluée par du permanganate de potassium à \qty{2,00e-2}{mol.L^{-1}}, en milieu acide :
$V_\mathrm{E}=\qty{17,6}{mL}$.
\begin{donnees}
Couples $\ce{MnO4-}/\ce{Mn^{2+}}$ et $\ce{O2}/\ce{H2O2}$ ; volume molaire à \qty{0}{\degreeCelsius} : \qty{22,4}{L.mol^{-1}} ; $\ce{2H2O2 -> 2H2O + O2}$.
\end{donnees}
\begin{questions}
\item Établir l'équation de la réaction support.
\item Calculer la concentration en peroxyde d'hydrogène de l'eau oxygénée.
\item En déduire son titre en volumes.
\item Quel volume équivalent attendait-on pour une eau oxygénée à exactement 10 volumes ? Proposer une explication à l'écart.
\end{questions}
\end{exercice}
\end{blocLivre}
\rubrique{Objectif DS}
\begin{blocLivre}{ex-5-12}
\begin{exercice}[3]{La vitamine C d'un jus de citron}
\begin{docu}{1}{Un besoin quotidien}
L'Agence nationale de sécurité sanitaire (Anses) fixe à \qty{110}{mg} par jour la référence nutritionnelle en vitamine C d'un adulte.
Cette vitamine est fragile : elle est oxydée par le dioxygène de l'air.
\end{docu}
\begin{docu}{2}{Le protocole}
On prélève \qty{10,0}{mL} de jus de citron fraîchement pressé et filtré, on ajoute de l'empois d'amidon, puis on titre par une
solution de diiode à \qty{2,00e-3}{mol.L^{-1}}. On recommence avec le même jus, conservé deux jours au réfrigérateur dans un
verre non couvert.
\par\smallskip\noindent\hfill\begin{tabular}{lcc}
\toprule
Jus & frais & après deux jours \\
\midrule
$V_\mathrm{E}$ (mL) & 14,2 & 11,4 \\
\bottomrule
\end{tabular}\hfill\null\par\smallskip
\end{docu}
\begin{donnees}
$\ce{C6H8O6 + I2 -> C6H6O6 + 2I- + 2H+}$ ; $M(\ce{C6H8O6})=\qty{176,0}{g.mol^{-1}}$.
\end{donnees}
\begin{questions}
\item Décrire le changement de teinte observé à l'équivalence et l'expliquer.
\item Calculer la concentration en masse de vitamine C du jus frais.
\item Quel pourcentage de vitamine C le jus a-t-il perdu en deux jours ?
\item \initiative Quel volume de jus frais couvre la référence nutritionnelle quotidienne ?
\end{questions}
\end{exercice}
\end{blocLivre}
\begin{blocLivre}{ex-5-13}
\begin{exercice}[3]{Le dioxyde de soufre du vin}
\begin{docu}{1}{Un additif réglementé}
Le dioxyde de soufre $\ce{SO2}$ protège le vin de l'oxydation et des micro-organismes. Il peut provoquer des réactions chez
les personnes sensibles : les étiquettes portent la mention « contient des sulfites ». La réglementation européenne limite
sa concentration à \qty{150}{mg.L^{-1}} dans un vin rouge sec et à \qty{200}{mg.L^{-1}} dans un vin blanc ou rosé sec.
\end{docu}
\begin{docu}{2}{Le titrage}
On titre $V_1=\qty{25,0}{mL}$ d'un vin blanc sec, acidifié et additionné d'empois d'amidon, par une solution de diiode à
\qty{5,00e-3}{mol.L^{-1}}. La teinte bleue persiste à partir de $V_\mathrm{E}=\qty{11,8}{mL}$.
\end{docu}
\begin{donnees}
Couples $\ce{I2}/\ce{I-}$ et $\ce{SO4^{2-}}/\ce{SO2}$ ; $M(\ce{SO2})=\qty{64,1}{g.mol^{-1}}$.
\end{donnees}
\begin{questions}
\item Écrire la demi-équation du couple $\ce{SO4^{2-}}/\ce{SO2}$ et établir l'équation de la réaction support.
\item Calculer la concentration en masse de dioxyde de soufre du vin.
\item Ce vin respecte-t-il la réglementation ?
\item \initiative Au-delà de quel volume équivalent un vin blanc sec serait-il non conforme ? Le même résultat serait-il
acceptable pour un vin rouge sec ?
\end{questions}
\end{exercice}
\end{blocLivre}
\begin{blocLivre}{ex-5-14}
\begin{exercice}[3]{Un bain antiseptique au permanganate}
\begin{docu}{1}{La préparation}
Pour un bain antiseptique, on dissout un comprimé de \qty{500}{mg} de permanganate de potassium $\ce{KMnO4}$ dans \qty{4,0}{L}
d'eau tiède : c'est une dilution « au 1/8000 », soit \qty{1}{g} pour \qty{8}{L}.
\end{docu}
\begin{docu}{2}{Un contrôle}
On titre $V_1=\qty{20,0}{mL}$ du bain par une solution d'ions fer(II) à \qty{5,00e-3}{mol.L^{-1}}, en milieu acide. La teinte violette
disparaît pour $V_\mathrm{E}=\qty{15,4}{mL}$.
\end{docu}
\begin{donnees}
$\ce{MnO4- + 8H+ + 5Fe^{2+} -> Mn^{2+} + 5Fe^{3+} + 4H2O}$ ; $M(\ce{KMnO4})=\qty{158,0}{g.mol^{-1}}$.
\end{donnees}
\begin{questions}
\item Quelle est ici l'espèce titrée ? le réactif titrant ? Expliquer le changement de teinte observé.
\item Écrire la relation à l'équivalence et calculer la concentration en ions permanganate du bain.
\item Calculer la concentration en masse de permanganate de potassium du bain.
\item \initiative Le bain a-t-il été correctement préparé ? Quel volume équivalent attendait-on ?
\end{questions}
\end{exercice}
\end{blocLivre}
\begin{blocLivre}{ex-5-15}
\begin{exercice}[3]{Lire un titrage sur un graphique}
\begin{docu}{1}{Le titrage}
On titre $V_1=\qty{20,0}{mL}$ d'une solution de diiode par du thiosulfate de sodium à \qty{5,00e-2}{mol.L^{-1}}, selon la réaction
$\ce{I2 + 2S2O3^{2-} -> 2I- + S4O6^{2-}}$. Une simulation donne l'évolution de trois espèces.
\schema{05-ex15}
\end{docu}
\begin{questions}
\item Associer chaque courbe à l'une des espèces $\ce{I2}$, $\ce{S2O3^{2-}}$ et $\ce{S4O6^{2-}}$, en justifiant.
\item Lire le volume équivalent et justifier ce choix.
\item \initiative Déterminer la concentration de la solution de diiode, de deux façons.
\item Décrire la teinte du mélange avant et après l'équivalence.
\end{questions}
\end{exercice}
\end{blocLivre}
